\documentclass[11pt,a4paper]{article}
\usepackage[utf8]{inputenc}
\usepackage[T1]{fontenc}
\usepackage[brazil]{babel}
\usepackage[a4paper,margin=2.5cm]{geometry}
\usepackage{amsmath,amssymb}
\usepackage{booktabs}
\usepackage{array}
\usepackage{calc}
\usepackage{graphicx}
\usepackage{hyperref}
\usepackage{parskip}
\usepackage{setspace}
\onehalfspacing
\title{}
\author{}
\date{}
\providecommand{\tightlist}{\setlength{\itemsep}{0pt}\setlength{\parskip}{0pt}}
\begin{document}
\maketitle
\subsection{Resumo}\label{resumo}

Este estudo piloto avalia empiricamente a capacidade de um ecossistema
de raciocínio determinístico (enumeração exata, aritmética de primos e
validação numérica) na resolução de quatro problemas canônicos do corpus
amostral IMO-AnswerBench (IMO Shortlist 2020--2022), usando um harness
de avaliação com escala 0--7 e \textbf{solver sem acesso à resposta
esperada} (anti-leak). Dos quatro problemas, dois foram resolvidos
exatamente por enumeração (teoria dos números: (p,q) = (7,3); álgebra: N
= 3, com esclarecimento semântico do termo ``quotient'' como divisão
inteira), um foi confirmado por enumeração com sensibilidade de
serialização na detecção da resposta (combinatória: 2\^{}\{n-1\}, n =
1..6), e um recebeu apenas evidência numérica forte, sem prova formal
(desigualdade algébrica). A taxa de acerto do harness foi 50\% (2/4),
com média 4.5/7. Identificamos três achados metodológicos: ambiguidade
semântica de paráfrases, divergência de definição de ``local maximum'' e
sensibilidade de serialização do grading. Os resultados são apresentados
comparados contextualmente à literatura (AlphaGeometry, AlphaGeometry2,
Aletheia), sem reivindicação de superioridade. Limitações: corpus
amostral pequeno (n = 4), sem head-to-head executado, solvers
determinísticos por domínio.

\subsection{1. Introdução}\label{introduuxe7uxe3o}

A resolução automatizada de problemas de olimpíadas de matemática
avançou significativamente com sistemas como AlphaGeometry e
AlphaGeometry2 (Busch et al., Nature 2024; versão 2025) e o agente de
pesquisa Aletheia da Google DeepMind, que relatam desempenho no nível de
medalha de ouro na IMO 2000--2024 (84\%). Esses resultados, entretanto,
dependem de busca neural, dedução simbólica especializada e
infraestrutura massiva. Este trabalho avalia um caminho complementar e
mais simples: \textbf{solvers determinísticos exatos} (sem aprendizado
de máquina) aplicados a um corpus amostral público do IMO-AnswerBench,
com protocolo anti-leak. O objetivo não é competir com sistemas neurais,
mas medir, de forma auditável e honesta, o que o raciocínio exato
alcança e onde falha --- contribuindo com um benchmark de linha de base
determinística.

\subsection{2. Método}\label{muxe9todo}

\subsubsection{2.1 Corpus}\label{corpus}

Quatro problemas canônicos do IMO-AnswerBench (Shortlist 2020--2022),
fixados no harness de avaliação: álgebra (N de uma soma de quocientes),
álgebra (desigualdade com parâmetro u), teoria dos números (equação
diofantina p\^{}3 - q\^{}5 = (p+q)\^{}2) e combinatória (permutações com
um único máximo local).

\subsubsection{2.2 Solver determinístico
(anti-leak)}\label{solver-determinuxedstico-anti-leak}

O solver recebe apenas o enunciado e nunca a resposta esperada. Para
cada categoria: enumeração exata de N (1..60) com divisão inteira;
enumeração exata de pares de primos (2..400); enumeração exata de
permutações (n = 1..6) com definição de máximo local incluindo extremos;
validação numérica extensiva da desigualdade para u = 2..6 e grade de t,
sem prova formal.

\subsubsection{2.3 Avaliação}\label{avaliauxe7uxe3o}

GradingHead na escala 0--7 (rubrica inspirada em benchmarking de provas
matemáticas). A resposta curta é detectada na solução gerada; a solução
é considerada correta quando a dedução exata atingiu a resposta e o
grading reconheceu a resposta.

\subsection{3. Resultados}\label{resultados}

\begin{longtable}[]{@{}
  >{\raggedright\arraybackslash}p{(\linewidth - 10\tabcolsep) * \real{0.1667}}
  >{\raggedright\arraybackslash}p{(\linewidth - 10\tabcolsep) * \real{0.1667}}
  >{\raggedright\arraybackslash}p{(\linewidth - 10\tabcolsep) * \real{0.1667}}
  >{\raggedright\arraybackslash}p{(\linewidth - 10\tabcolsep) * \real{0.1667}}
  >{\raggedright\arraybackslash}p{(\linewidth - 10\tabcolsep) * \real{0.1667}}
  >{\raggedright\arraybackslash}p{(\linewidth - 10\tabcolsep) * \real{0.1667}}@{}}
\toprule\noalign{}
\begin{minipage}[b]{\linewidth}\raggedright
Problema
\end{minipage} & \begin{minipage}[b]{\linewidth}\raggedright
Categoria
\end{minipage} & \begin{minipage}[b]{\linewidth}\raggedright
Resposta deduzida
\end{minipage} & \begin{minipage}[b]{\linewidth}\raggedright
Ground truth
\end{minipage} & \begin{minipage}[b]{\linewidth}\raggedright
Score (0--7)
\end{minipage} & \begin{minipage}[b]{\linewidth}\raggedright
Status
\end{minipage} \\
\midrule\noalign{}
\endhead
\bottomrule\noalign{}
\endlastfoot
algebra-001 & Álgebra & N = 3 & 3 & 6 & correto \\
algebra-004 & Álgebra & C = 2\^{}\{u-2\} (evidência numérica) &
2\^{}\{u-2\} & 3 & parcial \\
number-theory-001 & Teoria dos números & (7, 3) & (7, 3) & 6 &
correto \\
combinatorics-001 & Combinatória & 2\^{}\{n-1\} (n = 1..6) &
2\^{}\{n-1\} & 3 & confirmado por enumeração com ressalva de
serialização \\
\end{longtable}

Taxa de acerto (ground-truth reconhecido pelo grading): 2/4 (50\%).
Média: 4.5/7. Nenhuma solução obteve nota 7, pois nenhuma prova formal
completa foi produzida para todos os problemas.

\subsection{4. Achados metodológicos}\label{achados-metodoluxf3gicos}

\begin{enumerate}
\def\labelenumi{\arabic{enumi}.}
\tightlist
\item
  \textbf{Ambiguidade semântica (algebra-001)}: a paráfrase ``quotient
  of ab divided by N+1'' admite divisão racional (resultado N = 1) e
  divisão inteira (N = 3). A semântica oficial exige divisão inteira; o
  solver que a adota recupera o ground truth. Formalizações imprecisas
  alteram o resultado.
\item
  \textbf{Definição de ``local maximum'' (combinatorics-001)}: a
  definição restrita a interiores falha para n = 1, 2; a definição que
  inclui extremos confirma 2\^{}\{n-1\} para n = 1..6. A formalização da
  definição é parte do problema.
\item
  \textbf{Sensibilidade de serialização do grading}: a resposta deduzida
  ``2\^{}(n-1)'' em notação de programa não foi reconhecida pelo grading
  que espera ``2\^{}\{n-1\}'' (notação LaTeX). O reconhecimento de
  resposta curta é sensível à serialização --- limitação importante para
  benchmarks automáticos.
\end{enumerate}

\subsection{5. Discussão}\label{discussuxe3o}

Os resultados confirmam que raciocínio exato determinístico resolve
problemas de teoria dos números e álgebra de busca finita com total
auditabilidade, mas não substitui prova formal para desigualdades
paramétricas. A comparação com a literatura (AlphaGeometry/AG2/Aletheia)
é contextual e não competitiva: estes sistemas operam em escala muito
maior e com recursos neurais; nosso estudo estabelece uma linha de base
determinística reproduzível de 50\% num corpus amostral. A sensibilidade
do grading a serialização e ambiguidade de formalização recomenda que
benchmarks futuros publiquem especificações semânticas formais junto aos
enunciados.

\subsection{6. Conclusão}\label{conclusuxe3o}

Solvers determinísticos exatos, com anti-leak, atingem 50\% de acerto no
corpus amostral IMO-AnswerBench (2/4, média 4.5/7), resolvem
completamente problemas de enumeração (diofantina e soma de quocientes)
e fornecem evidência numérica forte --- não prova --- para desigualdade
paramétrica. Três achados metodológicos sobre ambiguidade, definição e
serialização são contribuições úteis para o desenho de benchmarks de
raciocínio matemático automatizado.

\subsection{7. Limitações}\label{limitauxe7uxf5es}

n = 4; sem head-to-head com sistemas externos; sem busca neural; sem
prova formal Lean4; solvers determinísticos por domínio; LLM on-device
auxiliou a estruturação textual, não os dados empíricos.

\subsection{Referências (contextuais)}\label{referuxeancias-contextuais}

\begin{itemize}
\tightlist
\item
  Trinh, T. H. et al.~AlphaGeometry. Nature 625, 468--475 (2024).
\item
  Google DeepMind. AlphaGeometry2 e Aletheia (comunicação de 2025).
\item
  Google DeepMind. IMO-AnswerBench / IMO-ProofBench / IMO-GradingBench
  (datasets abertos).
\end{itemize}
\end{document}
